\cleardoublepage

\section{实物系统部署与 Sim-to-Real 迁移}

\subsection{硬件平台架构}
为了验证 VLN-R1 在真实世界中的性能，我们搭建了一套基于移动机器人的具身智能实验平台。

\subsubsection{移动机器人底盘}
实验采用 TurtleBot4 为移动基座，其内部搭载了分布式控制系统。底层驱动通过 ROS2 Humble 框架实现对轮式升降、姿态传感器（IMU）以及激光雷达数据的管理。利用 Docker 容器技术，我们在底盘上实现了任务逻辑与系统驱动的解耦。

\subsubsection{Seeker 环视感知系统}
感知前端使用了 Seeker 系列 4-View 摄像头，能够提供覆盖约 $200^{\circ}$ 水平视场的高分辨率 RGB 图像与深度图。相比于传统的单目相机，该配置更接近仿真环境中的全景观测，为模型的跨域迁移提供了物理基础。

\subsection{边云协同的分布式推理框架}
考虑到多模态大模型对计算资源（显存、算力）的高要求，我们设计了一种“端-边-云”协同的分布式架构：
\begin{enumerate}
    \item **客户端 (Robot Side / Orin Nano)**：负责采集 Seeker 相机的 4 路视频流，进行去畸变与对齐预处理，并通过高带宽 TCP 协议发送至计算后端。
    \item **计算后端 (Cloud Side / RTX A6000)**：运行 VLN-R1 核心推理引擎。后端接收到观测数据后，通过逻辑推理（Think）生成下一步的目标位姿，并返回至物理终端。
    \item **控制执行 (Execution Layer)**：客户端接收到动作指令后，调用底层控制接口（如 \texttt{navigate\_to.py}）驱动机器人完成平稳移动。
\end{enumerate}

\subsection{Sim-to-Real 迁移策略}

\subsubsection{4-View 视场适配层}
仿真器中的 Panoramic 视角（12 帧）与实物 Seeker 的 4 视角存在分辨率和重叠度差异。我们在部署时引入了一个专门的视场适配层（FOV Adaption Layer），通过裁剪与重投影技术，使实车图像的特征分布尽可能贴近训练时的模拟数据。

\subsubsection{深度图归一化与增强}
实测深度图常含有噪声且有效距离有限（通常为 5 米）。为此，我们采用了两种策略：
\begin{itemize}
    \item **深度重标定**：将实测深度值映射至仿真环境的 $[0.1m, 10m]$ 区间。
    \item **AI 深度补全**：引入 FoundationStereo 模型对 Seeker 视差图进行后处理，填补了黑色反光物体带来的深度缺失。
\end{itemize}

\subsection{实地导航测试}
我们在真实实验室办公环境下进行了多次随机起始点的导航测试。实验结果表明，VLN-R1 能够准确识别真实场景中的“桌子”、“转角”、“房门”等路标。即便在光照变化剧烈或存在动态行人干扰的情况下，凭借其显式的思维链推理能力（如“检测到有人遮挡，我将原处等待”），系统展示出了极高的鲁棒性。实测成功率相比于传统的 End-to-End 方法有显著改善。
